% !TEX root = ../main.tex
\vspace{-10pt}
Il file \textit{fit\_utils.py} raccoglie il fit lineare pesato del package. L'obiettivo della release \texttt{1.0.0} è offrire un contratto semplice e stabile: una singola funzione pubblica, \texttt{lin\_fit}, e un tipo di ritorno immutabile, \texttt{LinearFitResult}.

\section{LinearFitResult}
\label{sec:linear-fit-result}
\begin{apidesc}
  \item[\textbf{Scopo}]
    Rappresentare in modo tipizzato e leggibile il risultato di un fit lineare.
  \item[\textbf{Campi pubblici}]
    \texttt{slope}, \texttt{intercept}, \texttt{slope\_std}, \texttt{intercept\_std},
    \texttt{covariance}, \texttt{correlation}, \texttt{residuals}, \texttt{residual\_std},
    \texttt{chi2}, \texttt{reduced\_chi2}, \texttt{dof}, \texttt{iterations},
    \texttt{converged}, \texttt{figure}.
  \item[\textbf{Note}]
    Il tipo è un dataclass immutabile; non sostituisce la funzione di fit ma formalizza il suo output.
\end{apidesc}

\section{\texorpdfstring{lin\_fit(x, y, sigma\_y, *, sigma\_x, xlabel, ylabel, title, decimals, tol, max\_iter, show\_plot, show\_band, show\_legend, show\_fit\_params, show\_grid, xlim, ylim, figsize, dpi, save\_path, title\_fontsize, title\_pad, legend\_fontsize, legend\_loc, data\_alpha, band\_alpha, grid\_alpha)}{lin\_fit(x, y, sigma_y, *, sigma_x, xlabel, ylabel, title, decimals, tol, max_iter, show_plot, show_band, show_legend, show_fit_params, show_grid, xlim, ylim, figsize, dpi, save_path, title_fontsize, title_pad, legend_fontsize, legend_loc, data_alpha, band_alpha, grid_alpha)}}
\label{sec:fit-lineare}
\begin{apidesc}
  \item[\textbf{Scopo}]
    Eseguire un fit lineare pesato \(y = mx + c\) su punti sperimentali con incertezze su \(y\) e, opzionalmente, su \(x\).
  \item[\textbf{Parametri chiave}]\hfill
    \begin{description}[font=\normalfont\ttfamily\bfseries, labelwidth=8em, leftmargin=8.5em, itemsep=1pt, parsep=0pt, topsep=2pt]
      \item[x, y] array-like $\to$ ascisse e ordinate dei dati
      \item[sigma\_y] array-like $\to$ incertezze su \(y\), strettamente positive
      \item[sigma\_x] array-like o \texttt{None} $\to$ incertezze su \(x\), strettamente positive
      \item[decimals] int (default \texttt{3}) $\to$ cifre decimali nella legenda del grafico
      \item[tol] float (default \texttt{1e-10}) $\to$ tolleranza relativa sulla convergenza iterativa
      \item[max\_iter] int (default \texttt{100}) $\to$ massimo numero di aggiornamenti dei pesi quando \texttt{sigma\_x} è attivo
      \item[show\_plot] bool (default \texttt{True}) $\to$ genera o meno la figura matplotlib
      \item[save\_path] stringa o \texttt{None} $\to$ salva la figura se \texttt{show\_plot=True}
    \end{description}
  \item[\textbf{Restituisce}]
    \texttt{LinearFitResult}.
  \item[\textbf{Eccezioni}]
    \texttt{ValueError} se i vettori non hanno la stessa lunghezza, contengono valori non finiti, hanno meno di 3 punti o contengono incertezze non positive;
    \texttt{ValueError} se \texttt{tol} o \texttt{max\_iter} sono invalidi;
    \texttt{ValueError} se \texttt{x} non contiene almeno due valori distinti;
    \texttt{ValueError} se \texttt{save\_path} è usato con \texttt{show\_plot=False};
    \texttt{RuntimeError} se il fit con \texttt{sigma\_x} non converge entro \texttt{max\_iter}.
\end{apidesc}

\paragraph{Fit con sole incertezze su \texttt{y}}
Nel caso base i pesi sono
\[
  w_i = \frac{1}{\sigma_{y_i}^2},
\]
e i coefficienti della retta sono stimati con
\[
  \hat{m} = \frac{\mathrm{Cov}_w(x,y)}{\mathrm{Var}_w(x)}, \qquad
  \hat{c} = \bar{y}_w - \hat{m}\,\bar{x}_w.
\]

\paragraph{Fit con incertezze anche su \texttt{x}}
Se viene passato \texttt{sigma\_x}, la funzione usa la varianza efficace
\[
  \sigma_{\mathrm{eff},i}^2 = \sigma_{y_i}^2 + \hat{m}^2 \sigma_{x_i}^2,
\]
aggiorna i pesi iterativamente e controlla la convergenza con \texttt{tol}. Il campo \texttt{iterations} conta quanti aggiornamenti dei pesi sono stati effettuati; il campo \texttt{converged} indica se il criterio di arresto è stato soddisfatto.

\paragraph{Diagnostiche}
Oltre ai parametri del fit, \texttt{LinearFitResult} contiene:
\begin{itemize}
  \item \texttt{residuals} e \texttt{residual\_std}
  \item \texttt{chi2} e \texttt{reduced\_chi2}
  \item \texttt{dof = n - 2}
  \item \texttt{figure}, che vale \texttt{None} se \texttt{show\_plot=False}
\end{itemize}

\paragraph{Esempio d'uso}
\begin{minted}{python}
fit = lin_fit(
    x,
    y,
    sigma_y,
    sigma_x=sigma_x,
    xlabel="tempo [s]",
    ylabel="spazio [m]",
    show_plot=False,
)

print(fit.slope)
print(fit.intercept_std)
print(fit.reduced_chi2)
\end{minted}
